\documentclass[11pt,a4paper]{article}

% --- pdfLaTeX / TeX Live 2024 safe preamble ---
\usepackage[T1]{fontenc}
\usepackage[utf8]{inputenc}
\usepackage{lmodern}
\usepackage{microtype}
\usepackage{amsmath,amssymb}
\usepackage[a4paper,margin=2.2cm]{geometry}
\usepackage{enumitem}
\usepackage[hidelinks]{hyperref}
\usepackage{graphicx}
\usepackage{float}
\usepackage{caption}

\title{Renal Physiology\\Chapters 40- Structured Rewrite}
\author{}
\date{}

\begin{document}
\maketitle

\noindent\textit{Source: Kam \& Power --- \textit{Principles of Physiology for the Anaesthetist}, 4th ed., Chapter 40-46.}

\tableofcontents

\section{Function of the Kidneys and Functional Anatomy}

\noindent\textit{Source: Kam \& Power --- \textit{Principles of Physiology for the Anaesthetist}, 4th ed., Chapter 40.}

\subsection{Overview}

The kidneys are central to \textbf{homeostasis}. Their main role is not simply urine production, but the continuous regulation of:
\begin{itemize}
    \item body water,
    \item electrolyte composition,
    \item acid--base status,
    \item excretion of metabolic waste,
    \item long-term blood volume,
    \item long-term arterial pressure.
\end{itemize}

This is achieved through three linked processes:
\begin{enumerate}
    \item \textbf{glomerular filtration},
    \item \textbf{tubular reabsorption},
    \item \textbf{tubular secretion}.
\end{enumerate}

A very large volume of plasma ultrafiltrate is formed each day, but most of it is reclaimed by the renal tubules. Final urine volume is therefore only a small fraction of filtered volume.

\subsection{Core Functions of the Kidneys}

\subsubsection{Water and electrolyte homeostasis}

This is the \textbf{primary function} of the kidneys.
\begin{itemize}
    \item The kidneys receive a high blood flow.
    \item About \textbf{180 L/day} of plasma ultrafiltrate is produced at the renal corpuscles.
    \item Most filtered water and solute are reabsorbed along the nephron.
    \item The \textbf{proximal tubule} alone reabsorbs about \textbf{60\% of filtered water and sodium}.
    \item Normal urine output is approximately \textbf{1.5 L/day}.
\end{itemize}

The kidneys therefore regulate:
\begin{itemize}
    \item extracellular fluid volume,
    \item osmolality,
    \item sodium balance,
    \item potassium balance,
    \item long-term blood volume,
    \item long-term arterial blood pressure.
\end{itemize}

\subsubsection{Excretion of waste products of metabolism}

The kidneys excrete metabolic waste products in urine, including:
\begin{itemize}
    \item \textbf{urea} --- from protein metabolism,
    \item \textbf{creatinine} --- from muscle metabolism,
    \item \textbf{uric acid} --- from nucleic acid metabolism,
    \item waste products derived from haemoglobin breakdown.
\end{itemize}

\subsubsection{Excretion of exogenous chemicals}

Many ingested substances are eliminated by the kidneys, including:
\begin{itemize}
    \item drugs,
    \item toxins,
    \item other exogenous chemicals.
\end{itemize}

Some substances are filtered, whereas others are actively secreted into the tubular lumen.

\subsubsection{Hormonal and metabolic functions}

The kidneys also function as an endocrine and metabolic organ. They produce:
\begin{itemize}
    \item \textbf{renin},
    \item \textbf{erythropoietin},
    \item the active form of vitamin D: \textbf{1,25-dihydroxyvitamin D$_3$}.
\end{itemize}

They also contribute to:
\begin{itemize}
    \item \textbf{gluconeogenesis} during starvation, especially from glutamine.
\end{itemize}

\subsubsection{Acid--base regulation}

By altering urinary excretion of:
\begin{itemize}
    \item hydrogen ions,
    \item bicarbonate,
\end{itemize}
the kidneys play a major role in long-term acid--base control.

\subsection{Fluid and Electrolyte Balance}

\subsubsection{Total body water}

Approximately \textbf{two-thirds of body weight} is water.

In a \textbf{70 kg adult male}, total body water is about \textbf{42 L}, divided into:
\begin{itemize}
    \item \textbf{intracellular fluid (ICF)}: \textbf{28 L},
    \item \textbf{extracellular fluid (ECF)}: \textbf{14 L},
    \begin{itemize}
        \item \textbf{plasma}: \textbf{3 L},
        \item \textbf{interstitial fluid}: \textbf{11 L}.
    \end{itemize}
\end{itemize}

Thus:
\begin{itemize}
    \item \textbf{ICF} is the largest compartment,
    \item \textbf{ECF = plasma + interstitial fluid}.
\end{itemize}

Total body water can be estimated by dilution methods using markers such as isotopically labelled water.

\subsubsection{Normal daily water balance}

Average daily water intake is around \textbf{2550 mL/day}.

This exceeds the obligatory losses required to maintain balance:
\begin{itemize}
    \item \textbf{urine}: \(\sim\)\textbf{430 mL} obligatory minimum,
    \item \textbf{insensible loss}: \(\sim\)\textbf{900 mL},
    \item \textbf{faeces}: \(\sim\)\textbf{100 mL},
    \item \textbf{sweat}: \(\sim\)\textbf{50 mL}.
\end{itemize}

Because intake usually exceeds obligatory loss, the kidneys normally excrete a relatively larger volume of dilute urine.

\subsubsection{Sodium balance}

Average daily sodium intake is approximately \textbf{457 mmol/day}.

Non-renal losses are small:
\begin{itemize}
    \item \textbf{sweat}: \(\sim\)\textbf{11 mmol/day},
    \item \textbf{faeces}: \(\sim\)\textbf{11 mmol/day}.
\end{itemize}

Therefore the kidneys must excrete the great majority, approximately \textbf{435 mmol/day}, to maintain sodium balance.

\subsubsection{Average minimum dietary requirements}

Approximate daily adult requirements are:
\begin{itemize}
    \item \textbf{water}: \textbf{25--35 mL/kg},
    \item \textbf{sodium}: \textbf{1--1.4 mmol/kg},
    \item \textbf{potassium}: \textbf{0.7--0.9 mmol/kg},
    \item \textbf{chloride}: \textbf{1.3--1.9 mmol/kg}.
\end{itemize}

As a rough practical approximation:
\begin{itemize}
    \item sodium requirement is about \textbf{100 mmol/day},
    \item potassium requirement is about \textbf{60 mmol/day}.
\end{itemize}

\subsection{Functional Anatomy of the Kidneys}

\subsubsection{The nephron}

The \textbf{nephron} is the functional unit of the kidney.

Each kidney contains about \textbf{1 million nephrons}.

A nephron consists of:
\begin{enumerate}
    \item \textbf{renal corpuscle},
    \item \textbf{proximal convoluted tubule},
    \item \textbf{loop of Henle},
    \item \textbf{distal convoluted tubule},
    \item \textbf{collecting ducts}.
\end{enumerate}

Flow of filtrate:
\[
\text{glomerular capillaries / Bowman's capsule}
\rightarrow
\text{proximal tubule}
\rightarrow
\text{loop of Henle}
\rightarrow
\text{distal convoluted tubule}
\rightarrow
\text{collecting duct}
\rightarrow
\text{papillary collecting ducts}
\rightarrow
\text{calyx}
\rightarrow
\text{renal pelvis}
\rightarrow
\text{ureter}
\]

The nephron wall is made of a \textbf{single layer of epithelial cells} separated from peritubular capillaries by a \textbf{basement membrane}.

In the distal nephron and collecting ducts, more than one epithelial cell type is present.

\subsubsection{Types of nephron}

Nephrons are classified according to the position of the renal corpuscle and the length of the loop of Henle.

\paragraph{Cortical nephrons}
\begin{itemize}
    \item Renal corpuscles lie more superficially in the cortex.
    \item They have \textbf{short loops of Henle}.
    \item Efferent arterioles form \textbf{peritubular capillaries} around nephron segments.
    \item Their main role is filtration and routine tubular processing.
\end{itemize}

\paragraph{Juxtamedullary nephrons}
\begin{itemize}
    \item Renal corpuscles lie close to the corticomedullary junction.
    \item They have \textbf{long loops of Henle} extending deep into the medulla.
    \item Efferent arterioles form both:
    \begin{itemize}
        \item \textbf{peritubular capillaries},
        \item \textbf{vasa recta}.
    \end{itemize}
    \item They are essential for \textbf{urine concentration and water conservation}.
\end{itemize}

\paragraph{Midcortical nephrons}
\begin{itemize}
    \item These may have either short or long loops of Henle.
\end{itemize}

\paragraph{Functional distinction}
\begin{itemize}
    \item \textbf{Cortical nephrons} are mainly associated with filtration.
    \item \textbf{Juxtamedullary nephrons} are especially important for concentrating urine.
\end{itemize}

\subsubsection{The renal corpuscle}

The renal corpuscle is located in the \textbf{renal cortex} and consists of:
\begin{itemize}
    \item a tuft of \textbf{glomerular capillaries},
    \item \textbf{Bowman's capsule}.
\end{itemize}

Its function is to form a \textbf{protein-free filtrate of plasma}.

Filtration occurs from glomerular capillaries into \textbf{Bowman's space} under the influence of opposing forces:
\begin{itemize}
    \item hydrostatic pressure favouring filtration,
    \item oncotic pressure opposing filtration.
\end{itemize}

The glomerulus is supplied by an \textbf{afferent arteriole} and drained by an \textbf{efferent arteriole}.

This arrangement is unusual because the capillary bed lies \textbf{between two arterioles}, allowing tight control of glomerular pressure and filtration.

\subsubsection{Epithelial specialisation along the nephron}

Different nephron segments are structurally specialised for different functions.

\paragraph{Proximal tubule}
The proximal tubule is the major site of bulk reabsorption. It reabsorbs large amounts of:
\begin{itemize}
    \item water,
    \item sodium,
    \item filtered solute.
\end{itemize}

This segment performs \textbf{isotonic reabsorption} and is responsible for reclaiming the majority of filtered fluid.

\paragraph{Loop of Henle}
The loop of Henle is important in creating the conditions required for \textbf{urine concentration}, especially in juxtamedullary nephrons.

\paragraph{Distal convoluted tubule and collecting duct}
These are more specialised regulatory segments.

Cell types include:
\begin{itemize}
    \item \textbf{principal cells},
    \item \textbf{type A intercalated cells},
    \item \textbf{type B intercalated cells}.
\end{itemize}

These segments fine-tune:
\begin{itemize}
    \item sodium handling,
    \item potassium handling,
    \item hydrogen ion secretion,
    \item bicarbonate handling,
    \item water reabsorption.
\end{itemize}

\subsection{How Kidney Structure Supports Kidney Function}

\subsubsection{High renal blood flow}

The kidneys receive a large blood flow, enabling:
\begin{itemize}
    \item rapid delivery of plasma for filtration,
    \item efficient waste clearance,
    \item close regulation of body fluid composition.
\end{itemize}

\subsubsection{Large filtration, small final urine volume}

A huge filtrate is produced each day, but most is reabsorbed.

This gives the kidney two major advantages:
\begin{itemize}
    \item effective elimination of waste products,
    \item precise control of final urine composition.
\end{itemize}

\subsubsection{Segmental processing along the nephron}

Each nephron segment has a distinct role:
\begin{itemize}
    \item \textbf{renal corpuscle}: filtration,
    \item \textbf{proximal tubule}: bulk reabsorption,
    \item \textbf{loop of Henle}: medullary concentration gradient,
    \item \textbf{distal nephron / collecting duct}: fine regulation.
\end{itemize}

\subsubsection{Juxtamedullary nephron design}

Long loops of Henle and associated vasa recta are essential for:
\begin{itemize}
    \item medullary hypertonicity,
    \item water conservation,
    \item formation of concentrated urine.
\end{itemize}

\subsection{Exam-Focused Synthesis}

\subsubsection{The kidneys do far more than produce urine}

They are central regulators of:
\begin{itemize}
    \item water balance,
    \item electrolyte balance,
    \item acid--base balance,
    \item blood volume,
    \item long-term arterial pressure.
\end{itemize}

\subsubsection{Three core renal processes}

All renal handling can be understood through:
\begin{enumerate}
    \item \textbf{filtration},
    \item \textbf{reabsorption},
    \item \textbf{secretion}.
\end{enumerate}

\subsubsection{Functional relevance of nephron type}

\begin{itemize}
    \item \textbf{Cortical nephrons}: predominantly filtration.
    \item \textbf{Juxtamedullary nephrons}: urine concentration.
\end{itemize}

\subsubsection{High-yield numerical facts}

\begin{itemize}
    \item Total body water \(\approx\) \textbf{two-thirds of body weight}.
    \item In a 70 kg male, total body water \(\approx\) \textbf{42 L}.
    \item Plasma volume \(\approx\) \textbf{3 L}.
    \item Interstitial fluid volume \(\approx\) \textbf{11 L}.
    \item Intracellular fluid volume \(\approx\) \textbf{28 L}.
    \item Plasma ultrafiltrate formed \(\approx\) \textbf{180 L/day}.
    \item Urine output \(\approx\) \textbf{1.5 L/day}.
    \item The proximal tubule reabsorbs \(\approx\) \textbf{60\% of filtered water and sodium}.
    \item Each kidney contains \(\approx\) \textbf{1 million nephrons}.
\end{itemize}

\subsection{Take-Home Points}

\begin{itemize}
    \item The kidneys are homeostatic organs whose main task is regulation of the internal environment.
    \item They maintain fluid, electrolyte, and acid--base balance while removing metabolic and exogenous waste.
    \item Their endocrine roles include production of renin, erythropoietin, and active vitamin D.
    \item The nephron is the structural and functional unit of the kidney.
    \item The renal corpuscle forms filtrate; the rest of the nephron modifies it by reabsorption and secretion.
    \item Cortical and juxtamedullary nephrons differ mainly in loop length and function.
    \item Juxtamedullary nephrons are essential for water conservation and concentrated urine production.
\end{itemize}

\section{Renal Blood Flow}

\subsection{Core facts}

The kidneys receive a disproportionately large blood flow relative to their size. Although they account for less than 1\% of body weight, they receive about 20\% of the cardiac output. This flow greatly exceeds metabolic requirements and is mainly necessary to generate a large volume of glomerular filtrate, allowing excretion of waste products.

Most renal blood flow is directed to the cortex, where the renal corpuscles are located. Medullary blood flow is much lower---about one-tenth of cortical flow---although still substantial in absolute terms. The principal resistance vessels in the renal circulation are the afferent and efferent arterioles; constriction of either reduces renal blood flow.

\subsection{Why renal blood flow is so high}

Renal blood flow is not primarily high because the kidney is metabolically demanding. Rather, the high flow is needed to support:
\begin{itemize}
    \item formation of a large glomerular filtrate
    \item regulation of fluid and electrolyte balance
    \item excretion of metabolic waste products
\end{itemize}

This is the key physiological reason the kidneys are perfused so heavily despite their small mass.

\subsection{Distribution of flow within the kidney}

\subsubsection{Cortex}

The majority of renal blood flow goes to the cortex. This supports filtration because the corpuscles are located there.

\subsubsection{Medulla}

The medulla receives much less flow, around 10-fold less than the cortex. Lower medullary flow is important because excessive medullary perfusion would tend to wash out the osmotic gradient needed for urinary concentration.

\subsection{Autoregulation of renal blood flow and GFR}

A major feature of the renal circulation is autoregulation.

\subsubsection{Definition}

Within a mean arterial pressure (MAP) range of about 75--170~mmHg, both:
\begin{itemize}
    \item renal blood flow (RBF), and
    \item glomerular filtration rate (GFR)
\end{itemize}
remain relatively constant. This is an intrinsic renal property and can be demonstrated even in an isolated kidney.

\subsubsection{Mechanism}

Autoregulation depends on changes in afferent arteriolar tone. The efferent arteriole is not involved in autoregulation as described here. When perfusion pressure rises or falls, the afferent arteriole adjusts its resistance correspondingly, stabilising both RBF and GFR.

\subsubsection{Importance}

Autoregulation prevents large pressure-dependent swings in:
\begin{itemize}
    \item renal perfusion
    \item filtration
    \item tubular load
    \item water and solute excretion
\end{itemize}

Without it, rises in arterial pressure would markedly increase renal blood flow and GFR, causing substantial solute and water loss.

\subsection{Mechanisms of autoregulation}

\subsubsection{Myogenic mechanism}

The smooth muscle in the afferent arteriole responds directly to changes in transmural pressure.
\begin{itemize}
    \item Increased perfusion pressure stretches the vessel wall
    \item stretch causes afferent arteriolar constriction
    \item resistance rises, limiting the increase in flow
\end{itemize}

Conversely, when perfusion pressure falls, resistance falls. This is the myogenic response.

\subsubsection{Tubuloglomerular feedback}

This is the second major autoregulatory mechanism.

\paragraph{Anatomical basis}

The macula densa, part of the juxtaglomerular apparatus, lies in the wall of the ascending limb of the loop of Henle close to the renal arterioles.

\paragraph{Effector substance}

The macula densa controls afferent arteriolar smooth muscle via adenosine, which acts as a vasoconstrictor.

\paragraph{Sequence when perfusion pressure rises}

\begin{enumerate}
    \item Renal perfusion pressure increases
    \item Glomerular capillary pressure rises
    \item Glomerular filtration increases
    \item More Na$^+$ and Cl$^-$ reach the ascending limb / macula densa region
    \item The macula densa releases more adenosine
    \item The afferent arteriole constricts
    \item Glomerular capillary pressure and GFR fall back toward normal
\end{enumerate}

\paragraph{Sequence when perfusion pressure falls}

The macula densa may produce nitric oxide as a vasodilator when renal perfusion pressure falls.

\paragraph{Functional role}

Tubuloglomerular feedback links tubular NaCl delivery to afferent arteriolar tone, thereby matching filtration to downstream tubular handling.

\subsection{Pressure diuresis}

Autoregulation is more precise for renal blood flow than for GFR over the MAP range of 75--170~mmHg. Therefore, as MAP rises:
\begin{itemize}
    \item filtration fraction increases
    \item glomerulotubular balance is disturbed
    \item urine flow increases
\end{itemize}

\subsubsection{Juxtamedullary nephron contribution}

Blood flow to juxtamedullary nephrons is not autoregulated. As blood pressure rises, blood is diverted from cortical nephrons to the juxtamedullary nephrons and vasa recta. This tends to wash out medullary solutes, specifically:
\begin{itemize}
    \item Na$^+$
    \item Cl$^-$
    \item urea
\end{itemize}

Loss of these solutes from the medulla reduces the kidney's ability to concentrate urine, increasing urinary loss of water and solutes. This is termed pressure diuresis.

\subsection{Sympathetic control of renal blood flow}

The kidneys receive noradrenergic sympathetic innervation. These nerves constrict both the afferent and efferent arterioles.

\subsubsection{Effect on renal blood flow}

Increased sympathetic activity reduces renal blood flow.

\subsubsection{Effect on GFR}

The effect on GFR is smaller than the effect on total renal blood flow.

This is because two opposing influences occur:
\begin{itemize}
    \item constriction of afferent and efferent arterioles increases glomerular capillary hydrostatic pressure, which tends to support filtration
    \item reduced renal blood flow increases glomerular capillary oncotic pressure, which opposes filtration
\end{itemize}

The net direct $\alpha$-adrenergic effect is a small reduction in GFR.

\subsubsection{Sympathetic activity and renin release}

Renal sympathetic nerves also stimulate renin secretion from the granular cells of the juxtaglomerular apparatus.

Renin release is increased by:
\begin{itemize}
    \item a direct $\beta_1$ effect of sympathetic nerve activity
    \item reduced GFR leading to reduced NaCl delivery to the macula densa, which further promotes renin release
\end{itemize}

This increases circulating angiotensin II.

\subsection{Angiotensin II}

Angiotensin II constricts both afferent and efferent arterioles, thereby reducing renal blood flow. Its haemodynamic effects on glomerular hydrostatic and oncotic pressures are similar to those of sympathetic activation.

However, angiotensin II has an additional action:
\begin{itemize}
    \item it causes mesangial cell contraction
    \item this reduces the filtration coefficient ($K_f$)
\end{itemize}

Because of this, the overall effect is a significant reduction in GFR.

\subsection{Eicosanoids}

Products of arachidonic acid metabolism affect renal resistance vessels.

\subsubsection{Vasodilator prostanoids}

The kidney produces locally active vasodilators including:
\begin{itemize}
    \item prostacyclin
    \item PGI$_2$
    \item PGE$_2$
\end{itemize}

These become important in clinical states where circulating vasoconstrictors are elevated.

\subsubsection{Clinical implication}

If prostaglandin synthesis is inhibited by cyclooxygenase inhibitors, renal blood flow may fall and renal function may deteriorate, particularly when the kidney is relying on these local vasodilators.

\subsection{Atrial natriuretic peptide (ANP)}

ANP increases glomerular hydrostatic pressure and GFR by:
\begin{itemize}
    \item dilating the afferent arteriole
    \item constricting the efferent arteriole
\end{itemize}

It is produced by atrial cells, and secretion is stimulated by atrial distension.

\subsection{Condensed exam summary}

\subsubsection{Numbers}

\begin{itemize}
    \item Kidneys: less than 1\% body weight
    \item Renal blood flow: about 20\% cardiac output
    \item Medullary flow: about 10 times less than cortical flow
    \item Autoregulation range: MAP 75--170~mmHg
\end{itemize}

\subsubsection{Main resistance vessels}

\begin{itemize}
    \item Afferent arteriole
    \item Efferent arteriole
\end{itemize}

\subsubsection{Autoregulation}

\begin{itemize}
    \item Keeps RBF and GFR relatively constant
    \item Depends on afferent arteriolar tone
    \item Mechanisms:
    \begin{itemize}
        \item myogenic
        \item tubuloglomerular feedback via macula densa / adenosine
    \end{itemize}
\end{itemize}

\subsubsection{Pressure diuresis}

\begin{itemize}
    \item Juxtamedullary flow not autoregulated
    \item Raised MAP increases vasa recta flow
    \item Medullary solute washout reduces concentrating ability
    \item Urinary water and solute loss rise
\end{itemize}

\subsubsection{Neurohumoral regulators}

\begin{itemize}
    \item Sympathetic nerves: $\downarrow$RBF, slight $\downarrow$GFR, $\uparrow$renin
    \item Angiotensin II: $\downarrow$RBF, significant $\downarrow$GFR, $\downarrow K_f$
    \item Prostaglandins / PGI$_2$ / PGE$_2$: renal vasodilation, protective in vasoconstrictor states
    \item ANP: afferent dilation + efferent constriction $\rightarrow$ $\uparrow$GFR
\end{itemize}

\subsection{Source}

Structured rewrite based on Chapter 41, ``Renal Blood Flow,'' in \textit{Principles of Physiology for the Anaesthetist} (Kam and Power). Source-constrained rewrite for exam revision.

\section{Glomerular filtration}

In the \textbf{renal corpuscle}, fluid passes from the \textbf{glomerular capillaries} into \textbf{Bowman’s space}. This movement is a \textbf{filtration process} driven by the balance of physical forces acting across a specialised filtration barrier.

The filtration barrier consists of three principal components:
\begin{enumerate}
    \item the \textbf{fenestrated endothelium} of the glomerular capillaries
    \item the \textbf{basement membrane}
    \item the \textbf{slit diaphragms between podocyte foot processes}
\end{enumerate}

Together, these structures permit rapid movement of water and small dissolved substances while restricting the passage of cells and most proteins.

The immediate driving forces for filtration are the opposing \textbf{hydrostatic} and \textbf{oncotic} pressures across this barrier. The fluid that enters Bowman’s space is therefore a \textbf{protein-free filtrate of plasma}.

A central determinant of filtration is the \textbf{glomerular capillary hydrostatic pressure}. This pressure depends on:
\begin{itemize}
    \item \textbf{renal blood flow}
    \item autoregulation of renal blood flow
    \item the resistance of the \textbf{afferent} and \textbf{efferent arterioles}
\end{itemize}

\subsection{Factors determining glomerular filtration in the renal corpuscle}

The determinants of glomerular filtration can be divided into two broad components.

\subsubsection{Filtration coefficient (\texorpdfstring{$K_f$}{Kf})}

The filtration coefficient represents:
\begin{itemize}
    \item the \textbf{surface area} available for filtration
    \item the \textbf{permeability} of the glomerular membrane
\end{itemize}

\subsubsection{Net filtration pressure}

Net filtration pressure is the net force tending to move fluid across the membrane.

\[
\text{Glomerular filtration} = K_f \times \text{net filtration pressure}
\]

Thus, glomerular filtration depends both on the properties of the filter itself and on the pressure gradient across it.

\subsection{Net filtration pressure}

Net filtration pressure is determined by the balance of one force promoting filtration and two forces opposing it.

\paragraph{Force promoting filtration}
\begin{itemize}
    \item \textbf{Glomerular capillary hydrostatic pressure}
\end{itemize}

\paragraph{Forces opposing filtration}
\begin{itemize}
    \item \textbf{Bowman’s capsule hydrostatic pressure}
    \item \textbf{Glomerular capillary oncotic pressure}
\end{itemize}

Therefore:
\[
\text{NFP} =
\text{glomerular capillary hydrostatic pressure}
-
\text{Bowman’s capsule hydrostatic pressure}
-
\text{glomerular capillary oncotic pressure}
\]

Oncotic pressure in Bowman’s space is ignored because the filtrate contains very little protein.

Physiologically:
\begin{itemize}
    \item glomerular capillary hydrostatic pressure pushes fluid \textbf{out} of the capillary
    \item glomerular capillary oncotic pressure pulls water \textbf{back into} the capillary
    \item Bowman’s capsule hydrostatic pressure mechanically opposes filtration
\end{itemize}

\subsection{The net filtration pressure along the glomerular capillary}

Net filtration pressure is not constant along the length of the glomerular capillary.

As blood passes from the \textbf{afferent} end to the \textbf{efferent} end:
\begin{itemize}
    \item hydrostatic pressure falls slightly because of vascular resistance
    \item oncotic pressure rises because protein-free fluid is filtered out, concentrating plasma proteins
\end{itemize}

Thus, the net filtration pressure falls along the capillary.

\subsubsection{At the afferent end}

Approximate values are:
\begin{itemize}
    \item glomerular capillary hydrostatic pressure = \textbf{60 mmHg}
    \item Bowman’s capsule hydrostatic pressure = \textbf{15 mmHg}
    \item glomerular capillary oncotic pressure = \textbf{21 mmHg}
\end{itemize}

\[
60 - 15 - 21 = 24 \text{ mmHg}
\]

So the net filtration pressure at the afferent end is about \textbf{24 mmHg}.

\subsubsection{At the efferent end}

Approximate values are:
\begin{itemize}
    \item glomerular capillary hydrostatic pressure = \textbf{58 mmHg}
    \item Bowman’s capsule hydrostatic pressure = \textbf{15 mmHg}
    \item glomerular capillary oncotic pressure = \textbf{33 mmHg}
\end{itemize}

\[
58 - 15 - 33 = 10 \text{ mmHg}
\]

So the net filtration pressure at the efferent end is about \textbf{10 mmHg}.

\subsubsection{Average net filtration pressure}

The average net filtration pressure across the glomerular capillary bed is approximately \textbf{17 mmHg}.

Despite this modest pressure, the glomerulus produces approximately \textbf{180 L/day} of filtrate. This is possible because the glomerular membrane has a very high filtration coefficient.

\subsection{Glomerular capillary hydrostatic pressure}

Glomerular capillary hydrostatic pressure is the principal \textbf{outward driving force} for filtration.

It tends to increase with rising arterial pressure, but this relationship is moderated by \textbf{autoregulation of renal blood flow}. As a result, changes in systemic arterial pressure do not translate directly into proportional changes in glomerular filtration under normal physiological conditions.

\subsubsection{Effect of afferent and efferent arteriolar tone}

The pressure within glomerular capillaries depends on the relative resistance of the vessels on either side.

\paragraph{Changes that increase glomerular capillary hydrostatic pressure}
\begin{itemize}
    \item \textbf{Afferent arteriolar dilatation}
    \item \textbf{Efferent arteriolar constriction}
\end{itemize}

These changes increase pressure within the glomerular capillary and therefore increase glomerular filtration. This is one effect of \textbf{atrial natriuretic peptide (ANP)}.

\paragraph{Changes that decrease glomerular capillary hydrostatic pressure}
\begin{itemize}
    \item \textbf{Afferent arteriolar constriction}
    \item \textbf{Efferent arteriolar dilatation}
\end{itemize}

These changes reduce glomerular capillary pressure and therefore reduce glomerular filtration.

\subsection{Bowman’s capsule hydrostatic pressure}

Bowman’s capsule hydrostatic pressure is an \textbf{opposing force} to filtration.

It reflects the pressure already present within Bowman’s space and the proximal urinary tract. If this pressure rises, filtration falls because the pressure gradient favouring movement of fluid into Bowman’s space is reduced.

The key applied point is that:
\begin{itemize}
    \item \textbf{obstruction to urine flow} raises Bowman’s capsule hydrostatic pressure
    \item this \textbf{reduces glomerular filtration}
\end{itemize}

This provides the physical explanation for reduced GFR in obstructive uropathy.

\subsection{Glomerular capillary oncotic pressure}

Glomerular capillary oncotic pressure is the second major force opposing filtration.

Because proteins are retained within the glomerular capillary while water is filtered out, protein concentration within capillary plasma rises as blood traverses the glomerulus. This progressively increases the oncotic force opposing further filtration.

The chapter highlights that:
\begin{itemize}
    \item reduced \textbf{renal blood flow}, for example from \textbf{renal sympathetic nerve activity}, can increase the rise in glomerular capillary oncotic pressure
    \item this higher oncotic pressure reduces net filtration pressure
\end{itemize}

Thus, reduced renal blood flow may lower filtration not only by reducing plasma delivery, but also by increasing the concentration effect of retained proteins.

\subsection{Filtration coefficient (\texorpdfstring{$K_f$}{Kf})}

The filtration coefficient is very high. It reflects:
\begin{itemize}
    \item the \textbf{permeability} of the glomerular capillary membrane
    \item the \textbf{surface area} available for filtration
\end{itemize}

Because \(K_f\) is high, an average net filtration pressure of only about \(17\ \text{mmHg}\) is sufficient to generate approximately \(180\ \text{L/day}\) of filtrate.

\subsubsection{Factors affecting \texorpdfstring{$K_f$}{Kf}}

\(K_f\) may be altered by factors affecting the dimensions of the \textbf{mesangial cells} situated between the glomerular capillaries. These cells can change the effective filtration surface area.

For example:
\begin{itemize}
    \item \textbf{angiotensin II} may reduce filtration by causing \textbf{mesangial cell contraction}
    \item this reduces glomerular surface area
    \item therefore \(K_f\) falls
    \item a lower \(K_f\) reduces glomerular filtration
\end{itemize}

This mechanism is distinct from changes in glomerular capillary hydrostatic pressure.

\subsection{The glomerular filtrate}

The filtrate entering Bowman’s space contains small solutes in concentrations similar to plasma.

The chapter specifies that the filtrate contains:
\begin{itemize}
    \item \textbf{electrolytes}
    \item \textbf{glucose}
    \item \textbf{amino acids}
\end{itemize}

These are present in the \textbf{same concentration as plasma}.

By contrast, the following are not normally filtered:
\begin{itemize}
    \item \textbf{cells}
    \item \textbf{large-molecular-weight molecules}
    \item \textbf{proteins}
\end{itemize}

Thus, the glomerular filtrate is best described as a \textbf{protein-free ultrafiltrate of plasma}.

\subsubsection{Selectivity by molecular size}

Molecular weight is an important determinant of filtration:
\begin{itemize}
    \item substances with molecular weight \textbf{below 7000 Da} are \textbf{freely filtered}
    \item above \textbf{7000 Da}, filtration progressively decreases
    \item \textbf{albumin}, with a molecular weight of about \textbf{70{,}000 Da}, is filtered only in very small quantities
\end{itemize}

\subsubsection{Selectivity by electrical charge}

The glomerular filtration barrier is coated with \textbf{negatively charged substances}. Therefore:
\begin{itemize}
    \item \textbf{negatively charged molecules} cross less easily
    \item plasma proteins are additionally restricted because they are also negatively charged
\end{itemize}

The barrier is therefore both:
\begin{itemize}
    \item \textbf{size-selective}
    \item \textbf{charge-selective}
\end{itemize}

\subsection{Filtration fraction}

Filtration fraction is the proportion of \textbf{renal plasma flow} that is filtered into the tubule.

The chapter gives:
\begin{itemize}
    \item renal plasma flow = \textbf{600 mL/min}
    \item GFR = \textbf{125 mL/min}
\end{itemize}

Therefore:
\[
\text{Filtration fraction} = \frac{125}{600} \approx 0.2
\]

So the filtration fraction is approximately \textbf{20\%}.

The rise in oncotic pressure within glomerular capillary plasma is directly proportional to the filtration fraction. The greater the proportion of plasma filtered, the greater the concentration of proteins remaining in the capillary, and therefore the greater the increase in glomerular capillary oncotic pressure.

\subsection{Chapter logic in one flow}

\begin{enumerate}
    \item Filtration occurs in the renal corpuscle across a specialised filtration barrier.
    \item The fluid entering Bowman’s space is a \textbf{protein-free filtrate of plasma}.
    \item The rate of filtration depends on:
    \begin{itemize}
        \item the \textbf{net filtration pressure}
        \item the \textbf{filtration coefficient}
    \end{itemize}
    \item Net filtration pressure is determined by:
    \begin{itemize}
        \item \textbf{glomerular capillary hydrostatic pressure} promoting filtration
        \item \textbf{Bowman’s capsule hydrostatic pressure} opposing filtration
        \item \textbf{glomerular capillary oncotic pressure} opposing filtration
    \end{itemize}
    \item Along the glomerular capillary:
    \begin{itemize}
        \item hydrostatic pressure falls slightly
        \item oncotic pressure rises
        \item net filtration pressure therefore falls
    \end{itemize}
    \item GFR is also influenced by \(K_f\), which depends on membrane permeability and filtration surface area.
    \item The filtrate contains \textbf{small solutes}, but not \textbf{cells} or significant \textbf{protein}.
    \item Filtration is restricted by both \textbf{molecular size} and \textbf{negative charge}.
    \item About \textbf{20\%} of renal plasma flow is filtered, giving a filtration fraction of about \textbf{0.2}.
\end{enumerate}

\subsection{Exam-oriented takeaways}

\subsubsection{Definitions}
\begin{itemize}
    \item \textbf{Glomerular filtration rate (GFR):} volume of fluid filtered from glomerular capillaries into Bowman’s space per unit time
    \item \textbf{Filtration coefficient (\(K_f\)):} product of permeability and filtration surface area
    \item \textbf{Net filtration pressure (NFP):} net force driving filtration
    \item \textbf{Filtration fraction:} GFR divided by renal plasma flow
\end{itemize}

\subsubsection{Key equations}
\[
\text{GFR} = K_f \times \text{NFP}
\]

\[
\text{NFP} =
\text{glomerular capillary hydrostatic pressure}
-
\text{Bowman’s capsule hydrostatic pressure}
-
\text{glomerular capillary oncotic pressure}
\]

\subsubsection{Normal values from the chapter}
\begin{itemize}
    \item afferent-end NFP = \textbf{24 mmHg}
    \item efferent-end NFP = \textbf{10 mmHg}
    \item mean NFP = \textbf{17 mmHg}
    \item GFR = \textbf{125 mL/min}
    \item filtrate formed = \textbf{180 L/day}
    \item renal plasma flow = \textbf{600 mL/min}
    \item filtration fraction = \textbf{20\%}
\end{itemize}

\subsubsection{Factors increasing GFR}
\begin{itemize}
    \item increased glomerular capillary hydrostatic pressure
    \item increased \(K_f\)
\end{itemize}

\subsubsection{Factors decreasing GFR}
\begin{itemize}
    \item reduced glomerular capillary hydrostatic pressure
    \item increased Bowman’s capsule hydrostatic pressure
    \item increased glomerular capillary oncotic pressure
    \item reduced \(K_f\)
\end{itemize}

\subsubsection{Specific examples named in the chapter}
\begin{itemize}
    \item \textbf{ANP}: increases filtration via afferent dilatation and efferent constriction
    \item \textbf{urinary obstruction}: reduces filtration by increasing Bowman’s capsule pressure
    \item \textbf{renal sympathetic activity}: reduces filtration partly by increasing glomerular capillary oncotic pressure
    \item \textbf{angiotensin II}: may reduce filtration by mesangial contraction and reduced \(K_f\)
\end{itemize}

\paragraph{Source.}
Kam and Power, \textit{Principles of Physiology for the Anaesthetist}, Chapter 42, ``Glomerular Filtration'', pp.\ 251--252.

\section{Tubular function}

\subsection{Tubular reabsorption and secretion}

The renal tubules reabsorb most of the water and solutes filtered at the glomerulus, while also secreting selected substances into the tubular lumen. The nephron epithelium consists of a single layer of cells separated from the peritubular capillaries by a basement membrane. Adjacent tubular cells are joined at their luminal borders by tight junctions.

Reabsorption from tubular lumen to blood occurs by two routes:
\begin{itemize}
    \item \textbf{Transcellular route}: across the luminal membrane, through the cell, and out across the basolateral membrane
    \item \textbf{Paracellular route}: between adjacent cells, across the tight junctions
\end{itemize}

Movement from the interstitial space into the peritubular capillaries is favoured because peritubular capillary hydrostatic pressure is lower than plasma oncotic pressure.

\subsection{Transport mechanisms in the renal tubule}

\subsubsection{Active transport}

Sodium is actively extruded from tubular cells across the basolateral membrane by the Na\textsuperscript{+}/K\textsuperscript{+}-ATPase. This lowers intracellular sodium concentration below that in the tubular fluid. Sodium can then move into the cell from the tubular lumen by several mechanisms. Sodium reabsorption therefore depends fundamentally on the basolateral Na\textsuperscript{+}/K\textsuperscript{+}-ATPase. This process also drives the reabsorption of glucose, amino acids, chloride, potassium and water.

\subsubsection{Facilitated diffusion}

Certain substances move down their electrochemical gradient by binding to specific membrane carrier proteins.

\subsubsection{Secondary active transport}

Two solutes may cross the membrane together via the same carrier. One moves down its electrochemical gradient and provides the energy for the other to move against its gradient.

\begin{itemize}
    \item \textbf{Co-transport}: both substances move in the same direction
    \item \textbf{Counter-transport}: the substances move in opposite directions
\end{itemize}

\subsection{Mechanism of tubular reabsorption}

\subsubsection{Sodium reabsorption}

\paragraph{Sodium reabsorption in the proximal tubule}

In the proximal tubule, sodium enters the cell across the luminal membrane:
\begin{itemize}
    \item by co-transport with glucose and amino acids
    \item by counter-transport with hydrogen ions or ammonium
\end{itemize}

These mechanisms depend on the basolateral Na\textsuperscript{+}/K\textsuperscript{+}-ATPase, which maintains both a low intracellular sodium concentration and a negative intracellular potential. Sodium therefore enters from the tubular lumen down its electrochemical gradient, and this movement supplies the energy for coupled transport of other solutes.

\paragraph{Glomerulotubular balance}

The proximal tubule reabsorbs sodium in proportion to the glomerular filtration rate (GFR). An increase in GFR is accompanied by an increase in proximal sodium reabsorption. This is known as \textbf{glomerulotubular balance}.

This occurs by two main mechanisms:
\begin{enumerate}
    \item \textbf{Increased filtered load of glucose and amino acids.} A rise in GFR increases the filtered load of glucose and amino acids. Since sodium reabsorption is linked to reabsorption of these substances, sodium and water reabsorption also increase.
    \item \textbf{Altered peritubular capillary Starling forces.} When GFR rises, protein concentration in glomerular capillary plasma increases. The blood entering the peritubular capillaries therefore has a higher oncotic pressure, which promotes movement of solute and water from the lateral intercellular space into the capillaries and enhances proximal reabsorption.
\end{enumerate}

The net effect is that a relatively constant fraction of filtered sodium and water is reabsorbed in the proximal tubule despite variations in GFR.

\paragraph{Sodium reabsorption in the thick ascending limb of the loop of Henle}

In the thick ascending limb, sodium enters the tubular cell from the lumen:
\begin{itemize}
    \item by co-transport with potassium and chloride
    \item by counter-transport with hydrogen ions
\end{itemize}

There is also significant paracellular sodium reabsorption in this segment. The lumen is electrically positive and the epithelium is relatively permeable to sodium, favouring passive movement by the paracellular route.

\paragraph{Sodium reabsorption in the distal convoluted tubule}

In the distal convoluted tubule, sodium enters the cell:
\begin{itemize}
    \item through specific sodium channels
    \item by co-transport with chloride
\end{itemize}

\paragraph{Sodium reabsorption in the collecting ducts}

In the principal cells of the collecting ducts, sodium enters the cell through luminal sodium channels.

\subsubsection{Water reabsorption in the nephron}

Water is reabsorbed by diffusion across cell membranes and tight junctions. Reabsorption of sodium and other solutes lowers luminal osmotic pressure, and water follows down the osmotic gradient.

The ascending limb of the loop of Henle is impermeable to water. Antidiuretic hormone (ADH) increases the water permeability of the collecting ducts.

\subsubsection{Chloride reabsorption in the nephron}

Chloride is reabsorbed by both paracellular and transcellular routes and, in general, follows sodium reabsorption.

In type B intercalated cells of the collecting duct, chloride reabsorption occurs independently of the Na\textsuperscript{+}/K\textsuperscript{+}-ATPase:
\begin{itemize}
    \item chloride enters across the luminal membrane by counter-transport with bicarbonate
    \item intracellular carbonic acid generates bicarbonate
    \item chloride leaves the cell through a basolateral chloride channel
\end{itemize}

\subsubsection{Glucose reabsorption in the nephron: transport maximum}

Glucose is reabsorbed in the proximal tubule by co-transport with sodium. Under normal conditions, all filtered glucose is reabsorbed.

However, the glucose transport system has a \textbf{transport maximum (Tm)}. When the filtered load exceeds this capacity, glucose appears in the urine.

In humans, with a normal GFR of 125~mL\,min\textsuperscript{-1}:
\begin{itemize}
    \item glucose begins to appear in urine when plasma glucose is about 10--12~mmol\,L\textsuperscript{-1}
    \item the transport system is fully saturated at about 15~mmol\,L\textsuperscript{-1}
\end{itemize}

Thus, glucosuria occurs when the filtered glucose load exceeds proximal tubular reabsorptive capacity, as in diabetes mellitus.

\subsubsection{Urea reabsorption in the nephron}

Approximately 50\% of filtered urea is reabsorbed in the proximal tubule by passive diffusion. As water is reabsorbed, luminal urea concentration rises and a diffusion gradient develops, favouring urea reabsorption.

The loop of Henle, distal convoluted tubule and cortical collecting ducts are relatively impermeable to urea, so luminal urea concentration rises further as water continues to be removed.

The inner medullary collecting duct is permeable to urea, and about a further 10\% of the filtered load is reabsorbed at this site. ADH increases urea permeability in the medullary collecting duct.

Overall, about 60\% of filtered urea is reabsorbed.

\subsubsection{Protein reabsorption in the nephron}

The glomerular membrane is relatively impermeable to large molecules, although a small amount of albumin is filtered. Most filtered protein is reabsorbed by the tubules, so normal urinary protein content is very low. Protein reabsorption has a transport maximum; if this is exceeded, protein appears in the urine.

\subsection{Mechanisms of tubular secretion}

\subsubsection{Proximal tubular secretion of organic anions}

The proximal tubule actively secretes organic anions into the lumen. Endogenous substances secreted in this way include:
\begin{itemize}
    \item urate
    \item bile salts
    \item fatty acids
    \item prostaglandins
\end{itemize}

Exogenous substances include:
\begin{itemize}
    \item para-aminohippuric acid (PAH)
    \item penicillin
    \item probenecid
    \item aspirin
\end{itemize}

Some organic anions, such as urate, are also filtered and reabsorbed, so secretion contributes to regulation of plasma concentration. Others, such as aspirin, are highly protein-bound and therefore poorly filtered, making proximal secretion important for their excretion.

\subsubsection{Proximal secretion of organic cations}

The proximal tubule also actively secretes organic cations, including:
\begin{itemize}
    \item creatinine
    \item acetylcholine
    \item catecholamines
    \item histamine
\end{itemize}

Drugs secreted by this pathway include:
\begin{itemize}
    \item pethidine
    \item morphine
    \item atropine
\end{itemize}

\subsubsection{Hydrogen ion secretion and bicarbonate reabsorption}

More than 4~mol of bicarbonate are filtered each day, and this alkali load must be reclaimed.

The mechanism depends on active hydrogen ion secretion into the tubular lumen. Hydrogen ions are generated in the tubular cell from intracellular carbonic acid formed from water and carbon dioxide. In the lumen, secreted hydrogen ions combine with filtered bicarbonate to form carbonic acid, which is then converted to carbon dioxide and water. Carbon dioxide diffuses into the cell, where intracellular carbonic anhydrase catalyses re-formation of carbonic acid and bicarbonate. The bicarbonate then passes into the peritubular capillaries.

Carbonic anhydrase is present:
\begin{itemize}
    \item inside proximal tubular cells
    \item on the luminal membrane of the proximal tubule
\end{itemize}

Hydrogen ion secretion occurs by three mechanisms:
\begin{enumerate}
    \item primary active H\textsuperscript{+}-ATPase
    \item Na\textsuperscript{+}/H\textsuperscript{+} counter-transport in the proximal tubule and ascending limb
    \item H\textsuperscript{+}/K\textsuperscript{+}-ATPase in type A intercalated cells of the collecting duct
\end{enumerate}

This process prevents loss of filtered bicarbonate but does not, by itself, acidify the urine. Urine becomes acidic when secreted hydrogen ions are excreted after buffering by urinary phosphate and other urinary buffers.

\subsection{Renal clearance}

\subsubsection{Definition}

For any substance, \textbf{renal clearance} is the volume of plasma completely cleared of that substance by the kidneys per unit time.

Units are expressed as volume per time, for example mL\,min\textsuperscript{-1} or L\,day\textsuperscript{-1}.

\subsubsection{Formula}

\[
\text{Renal clearance} = \frac{U \times V}{P}
\]

where:
\begin{itemize}
    \item $U$ = urine concentration of the substance
    \item $V$ = urine flow rate
    \item $P$ = plasma concentration of the substance
\end{itemize}

\subsubsection{Glucose}

Normally, glucose clearance is zero, because glucose is freely filtered but completely reabsorbed by the proximal tubule at normal plasma concentrations.

\subsubsection{Inulin}

Inulin is:
\begin{itemize}
    \item freely filtered
    \item not reabsorbed
    \item not secreted
    \item not metabolised
    \item not synthesised by the kidney
\end{itemize}

Therefore, inulin clearance is equal to GFR.

This makes inulin clearance a reliable measure of GFR, approximately 125~mL\,min\textsuperscript{-1} or 180~L\,day\textsuperscript{-1}.

\subsubsection{Para-aminohippuric acid}

PAH is filtered at the glomerulus and also secreted by the proximal tubule. At low plasma concentrations, almost all PAH delivered to functioning renal tissue is cleared in a single passage. Therefore:
\[
\text{PAH clearance} \approx \text{effective renal plasma flow}
\]

From this:
\[
\text{Effective renal blood flow} = \frac{\text{Effective renal plasma flow}}{1 - \text{haematocrit}}
\]

\subsubsection{Creatinine}

Creatinine is filtered and not reabsorbed, but a small quantity is secreted. Creatinine clearance therefore slightly overestimates true GFR.

Its limitations include:
\begin{itemize}
    \item plasma creatinine may not rise until approximately 50\% of renal function has been lost
    \item patients with low muscle mass may have a normal plasma creatinine despite substantial renal impairment
\end{itemize}

Age, weight and sex must therefore be considered. The text gives the Cockcroft--Gault formula:
\[
\text{Estimated GFR} = \frac{(140 - \text{age}) \times \text{weight (kg)}}{814 \times \text{plasma creatinine (mmol/L)}}
\]

\subsubsection{Urea}

Urea is freely filtered and 40--60\% is reabsorbed. Because tubular reabsorption varies, urea clearance is an unreliable measure of GFR.

\subsection{Summary of tubular handling of the glomerular filtrate}

\subsubsection{Proximal tubule}

Each day, 180~L of filtrate with an osmolality of 300~mOsm\,kg\textsuperscript{-1} H\textsubscript{2}O enters the proximal tubule.

The proximal tubule reabsorbs approximately:
\begin{itemize}
    \item 65\% of filtered sodium, chloride and water
    \item 55\% of filtered potassium
\end{itemize}

Active sodium reabsorption lowers luminal osmotic pressure, so water follows proportionately. This concentrates other luminal solutes such as potassium, chloride and urea, which are then reabsorbed by diffusion.

Much chloride reabsorption is paracellular, although active chloride transport also occurs. Sodium reabsorption in the proximal tubule also drives nutrient reabsorption and hydrogen ion secretion.

At the end of the proximal tubule:
\begin{itemize}
    \item 35\% of filtered water remains
    \item tubular fluid osmolality remains 300~mOsm\,kg\textsuperscript{-1} H\textsubscript{2}O
\end{itemize}

\subsubsection{Loop of Henle}

In the loop of Henle:
\begin{itemize}
    \item about 10\% of filtered water is reabsorbed in the descending limb
    \item about 25\% of filtered sodium and chloride is reabsorbed
    \item about 30\% of filtered potassium is reabsorbed
\end{itemize}

The loop generates the medullary osmotic gradient necessary for water reabsorption in the collecting ducts.

At the end of the loop:
\begin{itemize}
    \item 25\% of filtered water remains
    \item tubular fluid is hypo-osmotic, at about 100~mOsm\,kg\textsuperscript{-1} H\textsubscript{2}O
\end{itemize}

\subsubsection{Distal convoluted tubule}

Sodium and chloride reabsorption continue, but water is not reabsorbed, so the tubular fluid becomes more dilute.

At the end of the distal tubule:
\begin{itemize}
    \item 25\% of filtered water remains
\end{itemize}

\subsubsection{Collecting ducts}

ADH determines the water permeability of the collecting ducts.

\begin{itemize}
    \item With high ADH, less than 1\% of filtered water remains and urine osmolality may reach 1400~mOsm\,kg\textsuperscript{-1} H\textsubscript{2}O
    \item With no ADH, about 20\% of filtered water remains and urine output may rise to 30~L\,day\textsuperscript{-1}, with an osmolality below 100~mOsm\,kg\textsuperscript{-1} H\textsubscript{2}O
\end{itemize}

\subsubsection{Potassium handling in the distal convoluted tubule and collecting ducts}

\begin{itemize}
    \item Principal cells of the distal tubule and cortical collecting duct secrete potassium
    \item Type A intercalated cells reabsorb potassium
    \item The medullary collecting duct always reabsorbs potassium
\end{itemize}

With normal potassium intake, the distal nephron overall secretes potassium. During potassium depletion, the net effect becomes reabsorption.

\subsection{Hormonal control of tubular function}

\subsubsection{The renin--angiotensin system}

Aldosterone is produced by the adrenal cortex and promotes:
\begin{itemize}
    \item sodium reabsorption
    \item potassium secretion
\end{itemize}

Its principal site of action is the principal cell of the cortical collecting duct. Aldosterone secretion is closely controlled by the renin--angiotensin system.

\subsubsection{Renin}

Renin is secreted by the granular cells of the juxtaglomerular apparatus. It cleaves angiotensinogen, produced by the liver, to form angiotensin~I. Angiotensin-converting enzyme, particularly in pulmonary capillary endothelium, converts angiotensin~I to angiotensin~II.

Renin secretion is the rate-limiting step in this pathway.

\subsubsection{Control of renin secretion}

Renin secretion is increased by:
\begin{itemize}
    \item renal sympathetic nerves and circulating catecholamines acting via $\beta_{1}$ receptors
    \item intrarenal baroreceptors sensing reduced renal vascular pressure
    \item the macula densa when sodium chloride delivery falls
\end{itemize}

Reduced sodium chloride delivery to the macula densa may result from:
\begin{itemize}
    \item reduced GFR
    \item increased proximal tubular reabsorption
\end{itemize}

Renin secretion is directly inhibited by angiotensin~II through negative feedback.

\subsubsection{Angiotensin II}

Angiotensin~II reduces sodium and water loss and supports arterial pressure by:
\begin{itemize}
    \item stimulating sodium reabsorption indirectly through aldosterone
    \item directly stimulating tubular sodium reabsorption
    \item causing vasoconstriction and increasing sympathetic activity
    \item increasing peripheral resistance, cardiac output and arterial pressure
    \item constricting afferent and efferent arterioles, thereby reducing renal blood flow
    \item reducing the glomerular filtration coefficient ($K_f$) and significantly reducing GFR
    \item stimulating thirst and water intake
    \item stimulating ADH secretion
\end{itemize}

\subsubsection{Aldosterone}

Aldosterone release from the zona glomerulosa is stimulated by:
\begin{itemize}
    \item angiotensin~II
    \item increased plasma potassium
    \item adrenocorticotrophic hormone (ACTH)
\end{itemize}

In the cortical collecting duct, aldosterone:
\begin{itemize}
    \item increases sodium reabsorption
    \item increases potassium secretion
    \item increases hydrogen ion secretion by type A intercalated cells
\end{itemize}

It does so by increasing synthesis of proteins including:
\begin{itemize}
    \item the basolateral Na\textsuperscript{+}/K\textsuperscript{+}-ATPase
    \item luminal sodium channels
    \item luminal potassium channels
\end{itemize}

Aldosterone also increases sodium reabsorption from the gut, sweat glands and salivary glands.

\subsubsection{Antidiuretic hormone (ADH, vasopressin)}

ADH is synthesised in the hypothalamus and released from the posterior pituitary.

Its actions include:
\begin{itemize}
    \item increasing water permeability of the collecting duct by inserting water channels into the luminal membrane via cAMP
    \item increasing urea permeability of the inner medullary collecting duct
    \item causing vasoconstriction at high concentrations and potentially reducing renal blood flow and GFR
    \item stimulating sodium reabsorption and potassium secretion in the cortical collecting duct
\end{itemize}

\paragraph{Control of ADH secretion}

ADH release is regulated by osmoreceptors and baroreceptors acting through the hypothalamus.
\begin{itemize}
    \item increased plasma osmolality stimulates ADH secretion
    \item decreased plasma osmolality suppresses ADH secretion
\end{itemize}

The osmoreceptor set point lies between 280 and 295~mOsm\,kg\textsuperscript{-1} H\textsubscript{2}O.

A fall in plasma volume also stimulates ADH release, although this mechanism is less sensitive; typically a 5--10\% reduction in blood volume is required.

When plasma volume falls:
\begin{itemize}
    \item the ADH set point shifts to a lower osmolality
    \item the slope of the response becomes steeper
\end{itemize}

This permits water conservation even when plasma osmolality falls. The same receptor systems also influence thirst.

\subsubsection{Atrial natriuretic peptide (ANP) and brain natriuretic peptide}

Atrial stretch triggers release of ANP, which promotes sodium and water excretion by:
\begin{itemize}
    \item dilating the afferent arteriole
    \item constricting the efferent arteriole
    \item increasing net filtration pressure (NFP)
    \item increasing $K_f$
    \item inhibiting renin secretion
    \item inhibiting aldosterone release
    \item directly inhibiting sodium reabsorption in the collecting duct
\end{itemize}

A renal natriuretic peptide is also produced within the kidney. It increases GFR and reduces sodium reabsorption in the collecting ducts, but does not inhibit ADH action there.

\subsection{Control of renal sodium, water and potassium excretion}

\subsubsection{Control of renal sodium excretion}

Daily sodium filtered:
\[
140~\text{mmol\,L}^{-1} \times 180~\text{L\,day}^{-1} = 25\,200~\text{mmol\,day}^{-1}
\]

Most filtered sodium must therefore be reabsorbed. Urinary sodium excretion equals:
\[
\text{sodium filtered} - \text{sodium reabsorbed}
\]

Because plasma sodium concentration is normally maintained, renal sodium excretion depends on:
\begin{itemize}
    \item GFR
    \item tubular sodium reabsorption
\end{itemize}

Body sodium content determines extracellular fluid volume and therefore plasma volume and cardiovascular pressures. Changes in these pressures alter sodium excretion by changing both GFR and tubular reabsorption.

\paragraph{Filtration rate (GFR)}

As described in the preceding chapter:
\[
\text{Glomerular filtration} = K_f \times \text{NFP}
\]

\paragraph{Direct effects of body sodium on GFR}

When body sodium content falls and extracellular fluid volume contracts:
\begin{itemize}
    \item plasma volume falls
    \item arterial pressure falls
    \item glomerular capillary hydrostatic pressure falls
    \item glomerular capillary oncotic pressure rises
\end{itemize}

Therefore NFP falls and GFR decreases.

\paragraph{Indirect effects of body sodium on GFR}

When body sodium content falls:
\begin{itemize}
    \item baroreceptor reflexes increase renal sympathetic outflow
    \item renal blood flow falls
    \item GFR falls slightly
    \item activation of the renin--angiotensin system releases angiotensin~II, which markedly reduces GFR
    \item ANP production falls
\end{itemize}

These responses reduce sodium filtration and promote sodium conservation.

\paragraph{Body sodium, extracellular fluid volume and control of tubular sodium reabsorption}

Reduced body sodium increases tubular sodium reabsorption through:
\begin{itemize}
    \item aldosterone
    \item reduced renal interstitial hydrostatic pressure
    \item angiotensin~II
    \item renal sympathetic nerves
    \item reduced ANP
    \item increased ADH
\end{itemize}

These influences are reversed when extracellular volume is expanded.

\subsubsection{Control of renal water excretion}

Water excretion equals:
\[
\text{water filtered} - \text{water reabsorbed}
\]

The main control of water excretion is alteration of water reabsorption in the collecting ducts, which are water-permeable only in the presence of ADH.

Normally, mean daily urine output is about 1500~mL. In complete absence of ADH, as in diabetes insipidus, urine volume may rise to 25--30~L\,day\textsuperscript{-1}.

\subsubsection{Control of renal potassium excretion}

Potassium is the principal intracellular cation. Small changes in extracellular potassium concentration have major effects on membrane excitability, particularly in the heart.

If plasma potassium is 5~mmol\,L\textsuperscript{-1}, the daily filtered load is:
\[
5~\text{mmol\,L}^{-1} \times 180~\text{L\,day}^{-1} = 900~\text{mmol\,day}^{-1}
\]

Normally only 30--100~mmol\,day\textsuperscript{-1} is excreted, so most filtered potassium is reabsorbed. However, urinary potassium excretion can exceed the filtered load, demonstrating tubular secretion.

Urinary potassium excretion depends on:
\[
\text{filtered} + \text{secreted} - \text{reabsorbed}
\]

There is little regulation of potassium filtration or proximal and loop reabsorption:
\begin{itemize}
    \item the proximal tubule reabsorbs about 55\%
    \item the loop of Henle reabsorbs about 30\%
    \item the medullary collecting duct always reabsorbs potassium
\end{itemize}

Most regulation occurs in the principal cells of the distal tubule and cortical collecting duct, where potassium is secreted.

\paragraph{Control of tubular potassium secretion}

The main determinants are:
\begin{itemize}
    \item plasma potassium concentration
    \item aldosterone
\end{itemize}

Raised plasma potassium:
\begin{itemize}
    \item directly stimulates the basolateral Na\textsuperscript{+}/K\textsuperscript{+}-ATPase
    \item directly stimulates aldosterone secretion
\end{itemize}

Aldosterone increases synthesis of:
\begin{itemize}
    \item basolateral Na\textsuperscript{+}/K\textsuperscript{+}-ATPase
    \item luminal potassium channels
\end{itemize}

Potassium secretion also rises with increased tubular flow, because potassium that diffuses into the lumen is continuously washed away.

\paragraph{Effect of primary changes in body sodium balance on potassium excretion}

Primary changes in sodium balance do not have a major effect on potassium excretion.
\begin{itemize}
    \item Increased body sodium suppresses aldosterone but increases tubular flow
    \item Reduced body sodium raises aldosterone but decreases tubular flow
\end{itemize}

These opposing effects tend to balance each other.

\paragraph{Effect of primary changes in body water content on potassium excretion}

Primary alterations in body water content also do not greatly affect potassium excretion.
\begin{itemize}
    \item Increased body water lowers ADH but increases tubular flow
    \item Reduced body water raises ADH but decreases tubular flow
\end{itemize}

Again, these effects tend to offset one another.

\paragraph{Effect of alkalosis on potassium excretion}

Alkalosis increases urinary potassium excretion because a low plasma hydrogen ion concentration stimulates the basolateral Na\textsuperscript{+}/K\textsuperscript{+}-ATPase in principal cells.

\subsection{Mechanisms of action of diuretic drugs}

Mechanisms differ between drug classes.

\subsubsection{Thiazides}

Examples include bendroflumethiazide and chlorthalidone.

They reduce sodium reabsorption at the beginning of the distal convoluted tubule by inhibiting the luminal Na\textsuperscript{+}/Cl\textsuperscript{-} cotransporter.

\subsubsection{Loop diuretics}

Examples include furosemide, bumetanide and ethacrynic acid.

They reduce sodium reabsorption in the thick ascending limb by inhibiting the luminal Na\textsuperscript{+}/K\textsuperscript{+}/Cl\textsuperscript{-} cotransporter.

\subsubsection{Potassium-sparing diuretics}

There are two main subgroups.

\paragraph{Aldosterone antagonists}

Examples include spironolactone and potassium canrenoate.

These drugs antagonise aldosterone, mainly in the cortical collecting duct.

\paragraph{Amiloride and triamterene}

These are weak diuretics that promote potassium retention by blocking luminal sodium channels in the collecting ducts.

\subsubsection{Osmotic diuretics}

Mannitol is the main example.

It is filtered but not reabsorbed, so its osmotic effect within tubular fluid reduces water reabsorption.

\subsubsection{Carbonic anhydrase inhibitors}

Examples include acetazolamide and dichlorphenamide.

They reduce hydrogen ion secretion, causing bicarbonate and sodium reabsorption to fall. Their principal site of action is the proximal tubule.

\subsubsection{Mercurial diuretics}

An example is mersalyl.

These agents poison active transport pumps responsible for tubular ion reabsorption, including the Na\textsuperscript{+}/K\textsuperscript{+}-ATPase.

\subsubsection{Non-specific agents}

Agents that increase cardiac output may increase renal blood flow and urine output, including digoxin and plasma volume expanders. Low-dose dopamine acts as a renal vasodilator. Theophylline may increase GFR.

\subsubsection{Effect of diuretics on potassium excretion}

Most diuretics increase tubular flow to the cortical collecting duct and therefore increase potassium secretion. Consequently:
\begin{itemize}
    \item thiazides and loop diuretics may cause marked potassium depletion
    \item aldosterone antagonists, amiloride and triamterene do not increase urinary potassium loss in this way
\end{itemize}

\section{The Loop of Henle and Production of Concentrated Urine}

\subsection{Function of the loop of Henle}

To produce concentrated, hyperosmotic urine, water must be removed from tubular fluid while solute is retained. Water moves passively down an osmotic gradient, from a region of lower osmolality to one of higher osmolality. Therefore, if the kidney is to reabsorb water from tubular fluid, it must first generate a hyperosmotic environment outside the nephron. This is the role of the loop of Henle.

The loop of Henle creates a high medullary interstitial osmolality, which is essential for the formation of concentrated urine. Maximum urine osmolality is approximately 1200--1400~mOsm\,kg$^{-1}$ H$_2$O, whereas glomerular filtrate entering the nephron has an osmolality similar to plasma, about 300~mOsm\,kg$^{-1}$ H$_2$O. In the presence of ADH, water is reabsorbed from the collecting ducts because the surrounding medullary interstitium is strongly hyperosmotic.

\subsection{Counter-current arrangement}

The loop of Henle consists of two parallel limbs. Tubular fluid passes down into the medulla through the descending limb and then returns towards the cortex through the ascending limb. Flow in the two limbs is therefore in opposite directions, producing a counter-current arrangement.

\subsection{Descending limb of the loop of Henle}

The descending limb is permeable to water but impermeable to sodium and chloride. Consequently, when the medullary interstitium is hyperosmotic, water leaves the tubular lumen passively.

\subsection{Ascending limb of the loop of Henle}

In the thick ascending limb, sodium, potassium and chloride are actively reabsorbed by a co-transport mechanism. This transport of sodium together with potassium and chloride into the interstitium is the principal reason medullary interstitial osmolality becomes high.

The thick ascending limb is impermeable to water. There is also significant paracellular sodium reabsorption in this segment. This is favoured by the positive electrical potential in the tubular lumen together with the high sodium permeability of the epithelium.

\subsection{Effect on medullary interstitial fluid osmolality}

Active reabsorption of sodium and chloride from the thick ascending limb increases the osmolality of the surrounding interstitial fluid and simultaneously dilutes the fluid remaining within the tubular lumen. As interstitial osmolality rises, water leaves the descending limb.

Thus, within the loop of Henle, water reabsorption occurs in the descending limb, whereas sodium and chloride reabsorption occur in the ascending limb. The overall result is that both the interstitial fluid and the fluid within the descending limb become more concentrated, to around 400~mOsm\,kg$^{-1}$ H$_2$O, while fluid in the ascending limb becomes more dilute, to around 200~mOsm\,kg$^{-1}$ H$_2$O.

\subsection{Mechanism of counter-current multiplier effect}

Fluid entering the descending limb from the proximal tubule has an osmolality of about 300~mOsm\,kg$^{-1}$ H$_2$O, the same as plasma. Active transport of sodium and chloride from the ascending limb into the interstitial fluid raises interstitial osmolality. Water then leaves the descending limb until osmotic equilibrium is approached across that limb.

As fresh tubular fluid continues to flow, sodium and chloride are repeatedly pumped out of the ascending limb, and water repeatedly diffuses out of the descending limb. At each level of the loop, the difference in osmolality between the two limbs is maintained at about 200~mOsm\,kg$^{-1}$ H$_2$O. Although this transverse gradient is relatively small, repeated operation along the full length of the loop multiplies it into a large longitudinal corticomedullary gradient.

The result is a progressive increase in osmolality from the corticomedullary junction towards the tip of the loop of Henle, while the fluid leaving the ascending limb becomes progressively more dilute. With continued operation of this mechanism, together with the contribution of urea, the osmolality at the tip of the loop and in the deepest medullary interstitium rises to 1200--1400~mOsm\,kg$^{-1}$ H$_2$O, approximately 4--5 times plasma osmolality.

This mechanism is metabolically efficient because energy is required only to generate a local gradient of about 200~mOsm\,kg$^{-1}$ H$_2$O, while the arrangement of the loop amplifies this into a much larger longitudinal gradient. The process depends mainly on juxtamedullary nephrons, which comprise only about 15\% of all nephrons.

Tubular fluid leaving the ascending limb is hypo-osmotic relative to plasma, with an osmolality of about 100~mOsm\,kg$^{-1}$ H$_2$O. The distal tubule and cortical collecting duct remain impermeable to water unless ADH is present. When ADH is present, water diffuses out of these segments into the more concentrated interstitium.

\subsection{The role of urea}

NaCl contributes about half of the total medullary interstitial osmolality, around 600~mOsm\,kg$^{-1}$ H$_2$O. The remaining approximately 600~mOsm\,kg$^{-1}$ H$_2$O is contributed by urea.

Urea is freely filtered at the glomerulus, and about 50\% is reabsorbed in the proximal tubule. Therefore, the concentration of urea in the fluid entering the descending limb is only slightly greater than that in plasma. The loop of Henle and the cortical collecting duct are relatively impermeable to urea, so much of the filtered urea remains within the tubular fluid until it reaches the medullary collecting duct.

The medullary collecting duct is highly permeable to urea, and this permeability is increased further by ADH. By the time fluid reaches this segment, preceding water reabsorption has concentrated urea within the tubular lumen. Urea then diffuses down its concentration gradient from the medullary collecting duct into the interstitium, thereby making a major contribution to medullary interstitial osmolality.

Although the loop of Henle itself is relatively impermeable to urea, urea in the interstitial fluid surrounding the loop acts as an effective osmotic solute, drawing water from the descending limb and augmenting the concentrating effect produced by NaCl.

\subsection{The vasa recta: the counter-current exchanger}

The vasa recta are medullary capillary loops that run parallel and close to the descending and ascending limbs of the loop of Henle. They supply blood to the medulla without dissipating the osmotic gradient established there.

As blood descends into the medulla in the descending vasa recta, water leaves the vessel and solute, especially NaCl and urea, enters it. By the tip of the hairpin loop, blood osmolality may approach 1200~mOsm\,kg$^{-1}$ H$_2$O. As blood ascends back towards the cortex, the process reverses: water enters and solute leaves the ascending vasa recta. Consequently, blood leaving the medulla has an osmolality of only about 320~mOsm\,kg$^{-1}$ H$_2$O.

This counter-current exchange allows medullary perfusion while minimising washout of the osmotic gradient. In effect, the vasa recta trap solute within the medulla. The process is passive and is facilitated by the slow flow of blood through these vessels.

\subsection{The production of concentrated urine}

In the absence of ADH, the distal tubule and collecting ducts remain relatively impermeable to water. As a result, little water is reabsorbed from these segments, and the tubular fluid leaving them remains dilute, with an osmolality of less than 100~mOsm\,kg$^{-1}$ H$_2$O.

In the presence of ADH, the permeability of the collecting duct wall to water increases. Water is then reabsorbed from the cortical collecting duct until tubular fluid osmolality reaches about 300~mOsm\,kg$^{-1}$ H$_2$O, matching cortical interstitial osmolality. In the medullary collecting duct, because the surrounding interstitium is highly hyperosmotic, much more water is reabsorbed.

The final result is the excretion of a small volume of concentrated urine, with an osmolality of about 1400~mOsm\,kg$^{-1}$ H$_2$O, representing the obligatory minimum urine loss. About half of the osmolality of this concentrated urine is due to urea, while the remainder is contributed by sodium, chloride, potassium, creatinine and other solutes.


\section{Overview of Renal Control of Acid--Base Balance}

\subsection{Core principle}

The kidneys are essential for long-term acid--base homeostasis because they regulate plasma bicarbonate concentration and excrete the daily non-volatile acid load. On a normal diet, metabolism generates approximately 50--80 mmol of H$^+$ per day, and this acid must be excreted renally to prevent acidosis. This is physiologically important because the normal plasma hydrogen ion concentration is extremely low, around 36 nmol/L.

\subsection{Importance of the carbon dioxide--bicarbonate system}

The main extracellular buffer system is the carbon dioxide--bicarbonate buffer:
\[
CO_2 + H_2O \rightleftharpoons H^+ + HCO_3^-
\]

This buffer system is especially effective because it is an open physiological buffer:
\begin{itemize}
    \item the lungs regulate the acid component, $CO_2$
    \item the kidneys regulate the base component, $HCO_3^-$
\end{itemize}

Because each side of the buffer pair can be altered independently, the body can compensate for acid--base disturbances efficiently. In normal plasma:
\begin{itemize}
    \item $[CO_2] \approx 1.2$ mmol/L
    \item $[HCO_3^-] \approx 24$ mmol/L
    \item $pK_a = 6.1$
\end{itemize}

Using the Henderson--Hasselbalch equation:
\[
pH = pK_a + \log \frac{[HCO_3^-]}{[CO_2]}
\]

\[
pH = 6.1 + \log \frac{24}{1.2} = 7.4
\]

This explains why the kidneys can compensate for changes in PaCO$_2$ by changing plasma bicarbonate concentration.

\subsection{Main renal functions in acid--base balance}

The kidneys control acid--base status mainly by:
\begin{enumerate}
    \item reabsorbing filtered bicarbonate
    \item excreting hydrogen ions
    \item generating new bicarbonate
\end{enumerate}

Whenever H$^+$ is excreted in a form that is not simply used to reclaim filtered bicarbonate, this is equivalent to adding a new bicarbonate ion to blood.

\subsection{Relationship between hydrogen excretion and bicarbonate gain}

A central concept is:
\begin{itemize}
    \item excretion of H$^+$ in urine is equivalent to addition of HCO$_3^-$ to blood
    \item loss of HCO$_3^-$ in urine is equivalent to addition of H$^+$ to blood
\end{itemize}

This follows from the equilibrium:
\[
CO_2 + H_2O \rightleftharpoons H^+ + HCO_3^-
\]

If the kidney removes H$^+$ from the body, the reaction is driven such that bicarbonate is effectively added to plasma. Conversely, urinary bicarbonate loss is equivalent to acid retention.

\subsection{Why filtered hydrogen ions are negligible}

Although the kidney filters plasma continuously, the amount of filtered H$^+$ is tiny because plasma hydrogen ion concentration is extremely low.

Daily filtered H$^+$:
\[
36 \text{ nmol/L} \times 180 \text{ L/day} = 6840 \text{ nmol/day}
\]

This equals:
\[
0.684 \text{ mmol/day}
\]

This is negligible compared with the renal requirement to excrete approximately 40--80 mmol/day of acid. Therefore, acid excretion depends mainly on tubular secretion of H$^+$ rather than filtration.

\subsection{Why bicarbonate reabsorption is essential}

Daily filtered bicarbonate:
\[
24 \text{ mmol/L} \times 180 \text{ L/day} = 4320 \text{ mmol/day}
\]

Thus, more than 4 mol of bicarbonate are filtered each day. This large quantity cannot be lost in urine and must therefore be almost completely reabsorbed.

\subsection{Sites of bicarbonate reabsorption}

Filtered bicarbonate is reabsorbed along the nephron approximately as follows:
\begin{itemize}
    \item proximal tubule: 85\%
    \item thick ascending limb of the loop of Henle: 10\%
    \item distal convoluted tubule and cortical collecting duct: remainder
\end{itemize}

\subsection{Mechanism of bicarbonate reabsorption}

Filtered bicarbonate is not simply transported directly from lumen to blood. Reabsorption depends on tubular H$^+$ secretion.

\paragraph{Stepwise mechanism}
\begin{enumerate}
    \item Tubular cells secrete H$^+$ into the tubular lumen.
    \item Luminal H$^+$ combines with filtered HCO$_3^-$.
    \item This forms CO$_2$ and H$_2$O.
    \item CO$_2$ and H$_2$O diffuse into the tubular cell.
    \item Inside the cell, carbonic anhydrase catalyses re-formation of:
    \begin{itemize}
        \item HCO$_3^-$, which moves into peritubular blood
        \item H$^+$, which can be secreted again
    \end{itemize}
\end{enumerate}

\paragraph{Net effect}
\begin{itemize}
    \item one filtered bicarbonate ion is effectively returned to blood
    \item no net urinary H$^+$ excretion occurs in this step
\end{itemize}

Thus, this process represents reclamation of filtered bicarbonate rather than generation of new bicarbonate.

\subsection{Mechanisms of tubular hydrogen secretion}

All tubular segments possess a luminal H$^+$-ATPase.

Additional mechanisms vary by site:
\begin{itemize}
    \item proximal tubule and thick ascending limb: H$^+$ secretion also occurs via Na$^+$/H$^+$ counter-transport
    \item type A intercalated cells in the collecting duct: also possess H$^+$/K$^+$-ATPase, which secretes H$^+$ and reabsorbs K$^+$
\end{itemize}

\subsection{Bicarbonate secretion}

Type B intercalated cells in the cortical collecting duct can secrete bicarbonate into urine. The significance of this is not clearly established in the chapter, because the renal tubule normally reabsorbs most of the filtered bicarbonate.

\subsection{Generation of new bicarbonate}

New bicarbonate is added to blood when secreted H$^+$ is excreted with urinary buffers that are not reabsorbed.

\subsubsection{Phosphate mechanism}

If secreted H$^+$ combines with filtered phosphate in the tubular lumen, that H$^+$ is excreted in urine and a new HCO$_3^-$ is added to blood.

This pathway can excrete approximately 36 mmol H$^+$ per day.

\subsubsection{Ammonium and ammonia mechanism}

This is the more adaptable and higher-capacity mechanism.

\paragraph{Process}
\begin{enumerate}
    \item The liver produces glutamine from amino acid metabolism.
    \item Proximal tubular cells take up glutamine.
    \item They metabolise it to NH$_4^+$.
    \item NH$_4^+$ is secreted into the tubular lumen in exchange with Na$^+$.
    \item At the same time, new HCO$_3^-$ enters peritubular blood.
\end{enumerate}

Some NH$_3$ also diffuses from the medullary interstitium into the collecting duct, where it is protonated in tubular fluid, contributing to ammonium recycling.

The ammonium system has greater capacity for H$^+$ excretion than phosphate. The proximal tubule is the main site of glutamine metabolism and ammonium secretion.

\subsection{Minimum urinary pH}

The minimum urinary pH is about 4.4. Active hydrogen secretion becomes inhibited at very high luminal hydrogen ion concentrations.

\subsection{Fate of secreted hydrogen ions along the nephron}

The fate of secreted H$^+$ depends on which buffers remain in the tubular lumen.

\begin{itemize}
    \item In the proximal tubule and thick ascending limb, most secreted H$^+$ combines with bicarbonate.
    \item In the distal convoluted tubule and collecting ducts, little bicarbonate remains, so secreted H$^+$ combines mainly with phosphate.
\end{itemize}

\subsection{Control of renal hydrogen secretion}

Tubular H$^+$ secretion is:
\begin{itemize}
    \item increased by raised arterial CO$_2$ tension
    \item reduced by decreased arterial CO$_2$ tension
    \item increased by high extracellular H$^+$ concentration
    \item decreased when extracellular H$^+$ concentration is low
\end{itemize}

\subsection{Control of ammonium excretion}

Renal glutamine metabolism and NH$_4^+$ secretion are controlled by extracellular hydrogen ion concentration:
\begin{itemize}
    \item stimulated by acidosis
    \item inhibited by alkalosis
\end{itemize}

\subsection{Renal response to respiratory acid--base disorders}

\subsubsection{Respiratory acidosis}

When PaCO$_2$ rises:
\begin{itemize}
    \item extracellular H$^+$ concentration rises
    \item renal H$^+$ secretion increases
    \item NH$_4^+$ secretion increases
    \item more bicarbonate is reabsorbed and generated
\end{itemize}

The result is an increase in plasma bicarbonate concentration.

\subsubsection{Respiratory alkalosis}

When PaCO$_2$ falls:
\begin{itemize}
    \item renal H$^+$ secretion is inhibited
    \item NH$_4^+$ secretion is inhibited
    \item less bicarbonate is reabsorbed
\end{itemize}

The result is a reduction in plasma bicarbonate concentration.

\subsection{Renal response to metabolic acidosis}

In metabolic acidosis:
\begin{itemize}
    \item filtered bicarbonate load falls because plasma bicarbonate is low
    \item tubular H$^+$ secretion is sufficient to reabsorb all filtered bicarbonate
    \item H$^+$ is excreted as phosphate
    \item ammonium secretion increases because extracellular H$^+$ concentration is high
\end{itemize}

Therefore, the kidneys raise plasma bicarbonate by:
\begin{enumerate}
    \item reclaiming all filtered bicarbonate
    \item generating new bicarbonate through urinary H$^+$ excretion as phosphate and NH$_4^+$
\end{enumerate}

\subsection{Renal response to metabolic alkalosis}

In metabolic alkalosis:
\begin{itemize}
    \item filtered bicarbonate load rises because plasma bicarbonate is high
    \item tubular H$^+$ secretion is insufficient to reabsorb all filtered bicarbonate
    \item ammonium secretion decreases because extracellular H$^+$ concentration is low
\end{itemize}

Therefore, the kidneys lower plasma bicarbonate by allowing more filtered bicarbonate to remain in urine, that is, by not reabsorbing all filtered bicarbonate.

\subsection{High-yield numbers}

\begin{itemize}
    \item Daily non-volatile acid production: 50--80 mmol/day
    \item Plasma H$^+$ concentration: 36 nmol/L
    \item Plasma HCO$_3^-$ concentration: 24 mmol/L
    \item Plasma CO$_2$ concentration: 1.2 mmol/L
    \item Filtered bicarbonate per day: 4320 mmol/day
    \item Filtered hydrogen ions per day: 0.684 mmol/day
    \item Bicarbonate reabsorption:
    \begin{itemize}
        \item proximal tubule: 85\%
        \item thick ascending limb: 10\%
        \item remainder distally
    \end{itemize}
    \item Phosphate buffering capacity: approximately 36 mmol H$^+$/day
    \item Minimum urinary pH: 4.4
\end{itemize}

\section{Applied Physiology: Renal Failure}

\subsection{Definition and major categories}

Renal failure is broadly divided into acute kidney injury (AKI) and chronic kidney disease (CKD).

\paragraph{Acute kidney injury}
AKI is a reduction in kidney function developing over days. At the severe end of the spectrum, renal function may fall abruptly to the point that renal replacement therapy is required.

\paragraph{Chronic kidney disease}
CKD is a progressive decline in kidney function occurring over months to years.\cite{KamPowerRenalFailure}

\subsection{Aetiological classification}

Renal failure can be classified physiologically into prerenal, intrinsic renal, and postrenal causes.

\subsubsection{Prerenal renal failure}

Prerenal renal failure results from reduced renal perfusion rather than primary renal parenchymal disease. A typical example is hypovolaemia, in which reduced mean arterial pressure leads to compensatory afferent arteriolar vasoconstriction and a reduction in glomerular filtration.

Typical urine findings in prerenal oliguria reflect an intact kidney attempting to conserve salt and water:
\begin{itemize}
  \item high urine osmolality
  \item high urine urea concentration
  \item low urine sodium concentration.
\end{itemize}

\subsubsection{Intrinsic renal failure}

Intrinsic renal failure is caused by pathology within the kidney itself, including:
\begin{itemize}
  \item acute tubular necrosis
  \item toxic nephropathy
  \item inflammatory nephropathy
  \item autoimmune nephropathy.
\end{itemize}

\subsubsection{Postrenal renal failure}

Postrenal renal failure is caused by obstruction to urinary outflow. Back-pressure increases the pressure within Bowman’s capsule, reducing effective filtration pressure across the glomerulus.\cite{KamPowerRenalFailure}

\subsection{Acute versus chronic renal failure}

Acute and chronic renal failure can be distinguished using clinical and physiological features.

\begin{center}
\begin{tabular}{p{0.28\textwidth}p{0.28\textwidth}p{0.28\textwidth}}
\hline
\textbf{Feature} & \textbf{Acute} & \textbf{Chronic} \\
\hline
History & Short, days to weeks & Long, months to years \\
Haemoglobin & Usually normal & Reduced due to chronic anaemia \\
Renal size & Usually normal & Reduced \\
Renal osteodystrophy & Absent & Present \\
Peripheral neuropathy & Absent & Present \\
\hline
\end{tabular}
\end{center}

These differences reflect the longer-term endocrine, skeletal, and neurological consequences of chronic kidney disease.\cite{KamPowerRenalFailure}

\subsection{Core physiological consequences of renal failure}

The physiological effects of renal failure arise from impairment of three principal renal functions:
\begin{enumerate}
  \item regulation of fluid and electrolytes
  \item excretion of metabolic waste
  \item endocrine activity.
\end{enumerate}

In acute renal failure, the dominant problem is retention of water, electrolytes, and waste products. In chronic renal failure, endocrine failure becomes increasingly important.\cite{KamPowerRenalFailure}

\subsection{Failure of fluid and electrolyte balance}

\subsubsection{Initial compensatory adaptation}

As nephron number falls, surviving nephrons increase their workload by increasing filtration and tubular reabsorptive activity. This compensatory response is limited. Once approximately 75--80\% of nephrons are lost, sodium and water retention becomes clinically significant.\cite{KamPowerRenalFailure}

\subsubsection{Sodium and water retention}

When renal excretory capacity becomes insufficient:
\begin{itemize}
  \item sodium and water accumulate
  \item extracellular fluid volume expands
  \item total body water increases
  \item plasma osmolality tends to decrease.
\end{itemize}

Clinically, this may lead to dependent oedema, circulatory overload, and hypertension or other volume-related complications.\cite{KamPowerRenalFailure}

\subsubsection{Loss of concentrating and diluting ability}

Damage to renal architecture impairs the counter-current multiplier system. This is associated with:
\begin{itemize}
  \item loss of functioning juxtamedullary nephrons
  \item impaired generation of the medullary hyperosmotic gradient
  \item progressive loss of urinary concentrating ability.
\end{itemize}

Eventually, urine osmolality becomes fixed near plasma osmolality, producing isosthenuria. At this stage, the kidney can no longer appropriately vary urine concentration in response to changes in water balance.\cite{KamPowerRenalFailure}

\subsubsection{Handling of specific solutes}

Plasma sodium and chloride may remain near normal even with marked reduction in glomerular filtration rate (GFR), because tubular reabsorption of these ions decreases in parallel.

By contrast, once GFR falls below approximately 20--30\% of normal, plasma concentrations of the following rise disproportionately:
\begin{itemize}
  \item phosphate
  \item urate
  \item hydrogen ions.
\end{itemize}

\subsubsection{Metabolic acidosis}

Renal failure causes metabolic acidosis because the kidney can no longer:
\begin{itemize}
  \item excrete the daily fixed acid load effectively
  \item generate sufficient ammonia/ammonium buffering for secreted hydrogen ions.
\end{itemize}

As a consequence, bicarbonate falls and metabolic acids accumulate.\cite{KamPowerRenalFailure}

\subsubsection{Hyperkalaemia}

Potassium retention develops because the damaged kidney cannot adequately filter and secrete potassium. In acute kidney injury, the potassium load may be increased further by tissue injury, for example crush injury.

Hyperkalaemia is particularly dangerous because it impairs cardiac conduction and myocardial contractility, and may precipitate fatal arrhythmia or cardiac arrest.\cite{KamPowerRenalFailure}

\subsection{Failure of excretion of metabolic waste}

The kidney normally removes nitrogenous waste, particularly urea and creatinine, through glomerular filtration.

\paragraph{Creatinine}
Creatinine is filtered and essentially not reabsorbed, so serum creatinine rises as GFR falls.

\paragraph{Urea}
Urea also accumulates when filtration is reduced.

A key physiological point is that serum urea and creatinine usually do not rise above the normal range until GFR has fallen by approximately 50--60\%. A normal serum creatinine therefore does not exclude substantial loss of renal reserve.\cite{KamPowerRenalFailure}

\subsection{Failure of endocrine renal function}

Endocrine failure is much more prominent in chronic renal disease.

\subsubsection{Erythropoietin deficiency}

The kidney normally produces erythropoietin in peritubular cells. Loss of this function reduces red cell production and causes the anaemia of chronic renal failure.\cite{KamPowerRenalFailure}

\subsubsection{Renin--angiotensin--aldosterone system activation}

In early renal failure, ischaemic renal tissue may release more renin. This activates the renin--angiotensin--aldosterone system, causing:
\begin{itemize}
  \item angiotensin II mediated vasoconstriction
  \item aldosterone-mediated sodium and water retention.
\end{itemize}

This contributes to hypertension, particularly in earlier stages of renal impairment.\cite{KamPowerRenalFailure}

\subsubsection{Impaired vitamin D activation and mineral bone disease}

The kidney normally activates vitamin D. In renal failure:
\begin{itemize}
  \item vitamin D activation decreases
  \item intestinal calcium absorption decreases
  \item serum calcium falls
  \item serum phosphate rises because renal phosphate excretion is impaired.
\end{itemize}

Calcium and phosphate also complex together, reducing free calcium concentration further. This stimulates parathyroid hormone secretion.

The resulting pattern is:
\begin{itemize}
  \item hypocalcaemia
  \item hyperphosphataemia
  \item secondary hyperparathyroidism
  \item bone resorption and osteomalacia or renal bone disease.
\end{itemize}

\subsection{Major consequences of chronic renal failure}

The major consequences of chronic renal failure can be grouped into decreased excretion, decreased synthesis, and altered metabolism.

\subsubsection{Consequences of decreased excretion}

\begin{itemize}
  \item Uraemia or nitrogenous waste retention, leading to uraemic syndrome and platelet dysfunction
  \item Salt and water retention, leading to volume overload and hypertension
  \item Phosphate retention, leading to hyperparathyroidism and metastatic calcification
  \item Acid retention, leading to metabolic acidosis
  \item Potassium retention, leading to hyperkalaemia.
\end{itemize}

\subsubsection{Consequences of decreased synthesis}

\begin{itemize}
  \item Reduced erythropoietin production, causing anaemia
  \item Reduced vitamin D activation, causing osteomalacia and hyperparathyroidism.
\end{itemize}

\subsubsection{Consequences of altered metabolism}

\begin{itemize}
  \item Dyslipidaemia, contributing to accelerated atherosclerosis
  \item Altered sex hormone metabolism, contributing to abnormal reproductive function.
\end{itemize}

\subsection{Biochemical pattern in acute renal failure}

Typical biochemical abnormalities in acute renal failure are:
\begin{itemize}
  \item hyperkalaemia
  \item reduced bicarbonate
  \item raised urea
  \item raised creatinine
  \item raised uric acid
  \item hyperphosphataemia
  \item hypocalcaemia.
\end{itemize}

\subsection{Physiological distinction between prerenal and intrinsic renal dysfunction}

The chapter provides classic urine indices that help distinguish prerenal from intrinsic renal failure.

\subsubsection{Prerenal pattern}

This reflects preserved tubular function and appropriate conservation:
\begin{itemize}
  \item urine specific gravity high
  \item urine osmolality high
  \item urine sodium low
  \item fractional excretion of sodium low.
\end{itemize}

\subsubsection{Intrinsic renal pattern}

This reflects tubular damage and impaired conservation:
\begin{itemize}
  \item urine specific gravity low
  \item urine osmolality low
  \item urine sodium high
  \item fractional excretion of sodium high.
\end{itemize}

Conceptually, prerenal renal failure indicates that the kidney is responding appropriately to reduced perfusion, whereas intrinsic renal failure indicates that the kidney can no longer mount an appropriate tubular response.\cite{KamPowerRenalFailure}

\subsection{Investigations and assessment of renal function}

Investigations can be grouped into urinalysis, blood tests, and imaging or tissue diagnosis.

\subsubsection{Urinalysis}

Urinalysis includes:
\begin{itemize}
  \item specific gravity
  \item osmolality
  \item pH
  \item protein by dipstick
  \item microscopy for casts, white cells, and red cells.
\end{itemize}

\subsubsection{Blood and serum tests}

Relevant blood investigations include:
\begin{itemize}
  \item serum urea
  \item serum creatinine
  \item estimated GFR or creatinine clearance
  \item electrolytes, especially potassium, phosphate, calcium, and sodium
  \item arterial blood gases for assessment of acidosis
  \item cystatin C.
\end{itemize}

Novel biomarkers such as N-acetyl-glucosaminidase are mentioned as research tools.

\subsubsection{Imaging and tissue diagnosis}

The source lists:
\begin{itemize}
  \item plain abdominal X-ray
  \item excretion urography
  \item ultrasound
  \item CT
  \item MRI
  \item pyelography
  \item renal biopsy.
\end{itemize}

\end{document}
